\label{sec:self-consistency-anchoring-bias}

\textbf{System variability.} Across 10 repeated runs per patient, the System exhibited substantial output variability. Mean within-patient standard deviation was 0.192 (median 0.176), with mean score range of 0.554 (median 0.533). Only 14.2\% of patients (27/190) showed highly consistent outputs (standard deviation $<$0.1), while 5.3\% (10/190) showed very high variability (standard deviation $\geq$0.4). This indicates that the System's medication safety assessments are sensitive to sampling variability in the underlying language model, with most patients receiving meaningfully different evaluations across runs despite identical inputs.

\textbf{Scorer variability.} By contrast, the LLM-as-a-judge scorer demonstrated much higher consistency. Mean within-patient standard deviation was 0.041 (median 0.000), with 83.5\% of patients (167/200) showing minimal scorer variability (standard deviation $<$0.1). The scorer's mean standard deviation was 5.04-fold lower than the system's, showing that the evaluation methodology contributes minimal noise relative to genuine system output variability.

\textbf{Anchoring bias.} At the binary level, we calculated the conditional probability of flagging a patient given initial flagging status. For initially flagged cases, the probability of re-flagging on subsequent runs was 96.4\%, indicating high consistency in identifying genuinely problematic cases. For initially negative cases, the probability of remaining unflagged was reduced to 62.9\%, suggesting the system is more prone to inconsistency on cases where no intervention is needed, occasionally raising spurious concerns. Applying these self-consistency rates to the 277 evaluable patients yields expected true positives of 198.6 (206 × 0.964) and expected true negatives of 44.7 (71 × 0.629), producing a model self-consistency ceiling of 87.8\% accuracy. Observed clinician-agreement accuracy (95.7\%) exceeded this ceiling by 7.9\%.